\documentclass[10pt,a4paper]{amsart}
\usepackage[margin=19mm]{geometry}
\usepackage{amsmath,amssymb,mathtools}
\usepackage[scr=rsfs,cal=euler]{mathalfa}
\usepackage{xfrac}
\usepackage[unicode,hidelinks,bookmarks=false]{hyperref}
\newcommand{\ie}{i.e.\ }
\newcommand{\cf}{cf.\ }
\newcommand{\resp}{respectively\ }
\newcommand{\Subsection}{\subsection}
\providecommand{\nobreakdash}{\nobreak\hbox{-}}
\setlength{\emergencystretch}{2em}
\multlinegap0pt
\usepackage{mathtools}
\newtheorem{thm}{Theorem}[section]
\newtheorem{lem}[thm]{Lemma}
\newtheorem{prop}[thm]{Proposition}
\newtheorem{cor}[thm]{Corollary}
\newtheorem{defn}[thm]{Definition}
\newtheorem{example}[thm]{Example}
\newtheorem{notation}[thm]{Notation}
\newtheorem{remark}[thm]{Remark}
\newtheorem{question}[thm]{Question}
\newcommand{\R}{\mathbb{R}}
\newcommand{\N}{\mathbb{N}}
\newcommand{\ip}[2]{\langle #1,\, #2 \rangle}
\newcommand{\norm}[1]{\left\lVert #1 \right\rVert}
\newcommand{\bn}[2]{{\textstyle\binom{#1}{#2}}}
\DeclareMathOperator{\tr}{tr}
\DeclareMathOperator{\Tr}{Tr}
\DeclareMathOperator{\Sym}{Sym}
\DeclareMathOperator{\sgn}{sgn}
\DeclareMathOperator{\diag}{diag}
\DeclareMathOperator{\spn}{span}
\DeclareMathOperator{\ap}{ap}
\DeclareMathOperator{\PG}{PG}
\begin{document}
\title[Finite free convolution via reproducing kernels and squarefr]{Finite free convolution via reproducing kernels and squarefree algebras}
\author{Thomas Sinclair}
\date{}
\maketitle
\noindent\emph{Source and licence.} Finite free convolution via reproducing kernels and squarefree algebras. Version arXiv:2606.10870v2 (CC BY 4.0). This material is distributed under the Creative Commons Attribution 4.0 International licence, http://creativecommons.org/licenses/by/4.0/, as recorded in the arXiv OAI licence metadata for this exact version and shown on the abstract page https://arxiv.org/abs/2606.10870v2. Full rights notice: licenses/arxiv-2606.10870-CC-BY-4.0.txt.\par\smallskip
\noindent\textit{Editorial scope.} Counted unit: the original \texttt{thm} environment \texttt{thm:hamiltonian} at source lines 1023--1042 together with its complete proof at lines 1044--1046; the delivered document keeps source lines 929--1047 so that every referenced label and every citation resolves inside it. Native editable source is the paper's own concatenated TeX; the AMS wrapper is editorial formatting only. No publisher class, upstream script or unrelated artwork is redistributed.
\section{Trace--log expansion and examples}
\label{sec:laplacian}
%%%%%%%%%%%%%%%%%%%%%%%%%%%%%%%%%%%%%%%%%%%%%%%%%%%%%%%%%%%%%

We now apply the multilinear framework to determinantal examples.  In each case, the trace--log formula in $\mathcal{A}_n$ gives an explicit combinatorial description of the cumulants.

\subsection{The trace--log identity in $\mathcal{A}_n$}

Let $M\in M_n(\R)$ and set
\[
  D(t) := \diag(t_1,\dots,t_n),
  \qquad
  X := D(t)\, M \in M_n(\mathcal{A}_n).
\]
Since $t_i^2 = 0$ in $\mathcal{A}_n$, every entry of $X^m$ involves monomials of degree at least~$m$, so $X^{n+1} = 0$. Thus $X$ is nilpotent and $\log(I+X) = \sum_{m=1}^n \frac{(-1)^{m+1}}{m}\, X^m$ is well-defined.

\begin{prop}[Trace--log identity]
\label{prop:trace-log}
For every $M\in M_n(\R)$,
\[
  \log\det(I + D(t)\, M)
  = \Tr\log(I + D(t)\, M)
  = \sum_{m=1}^n \frac{(-1)^{m+1}}{m}\,
  \Tr\bigl((D(t)\, M)^m\bigr)
\]
in $\mathcal{A}_n$.
\end{prop}

\begin{proof}
The identity $\log\det(I+X) = \Tr\log(I+X)$ holds for nilpotent matrices over any commutative $\R$-algebra, as both sides agree as polynomial identities in the entries.
\end{proof}

\begin{prop}\label{prop:trace-cycle}
For each $m\ge 1$, the coefficient of $t^S$ (with $|S|=m$) in
$\Tr\bigl((D(t)\, M)^m\bigr)$ is
\[
  \sum_{\substack{(i_1,\dots,i_m)\\
  \{i_1,\dots,i_m\}=S}}
  M_{i_1 i_2}\, M_{i_2 i_3}\cdots M_{i_m i_1},
\]
where the sum is over all sequences $(i_1,\dots,i_m)$ that are
permutations of $S$.
\end{prop}

\begin{proof}
Expanding the matrix product and taking the trace, terms with
a repeated index vanish in $\mathcal{A}_n$ since $t_i^2=0$.
\end{proof}

Combining these gives the main formula.

\begin{thm}[Cycle expansion for multilinear cumulants]
\label{thm:cycle-expansion}
For $f_M(t) := \det(I+D(t)\, M)\in\mathcal{A}_n$ with
$\log f_M = \sum_{\varnothing\neq S} \kappa_S(M)\, t^S$,
\[
  \kappa_S(M)
  = (-1)^{k+1}
  \sum_{C(S)}
  M_{i_1 i_2}\, M_{i_2 i_3}\cdots M_{i_k i_1},
\]
where $k=|S|$ and $C(S)$ denotes classes of cyclic orderings of $S$ modulo rotation.
\end{thm}

\begin{proof}
The monomial $t^S$ arises only from the $m=|S|$ term.  Each
cyclic ordering is counted $k$ times by starting point, absorbing
the factor $1/k$.
\end{proof}

\begin{remark}\label{rem:small-cumulants}
For small sets, the formula gives:
\begin{itemize}
\item $|S|=1$:
$\kappa_{\{i\}}(M) = M_{ii}$.

\item $|S|=2$:
$\kappa_{\{i,j\}}(M) = -M_{ij}\, M_{ji}$.

\item $|S|=3$:
$\kappa_{\{i,j,k\}}(M) = M_{ij}\, M_{jk}\, M_{ki}
+ M_{ik}\, M_{kj}\, M_{ji}$.
\end{itemize}
The singleton cumulant records the diagonal entry, the pairwise
cumulant records (minus) the product of the off-diagonal pair,
and the triple cumulant records the two directed triangles
through the three vertices.
\end{remark}

\subsection{Graph Laplacians}

Let $G$ be a finite simple graph on $[n]$ with Laplacian $L_G$,
and set $f_G(t) := \det(I + D(t)\, L_G)$.


\begin{thm}[Hamiltonian-cycle formula]
\label{thm:hamiltonian}
Write $\log f_G = \sum_{\varnothing\neq S} \kappa_S(G)\, t^S$.
Then:
\begin{enumerate}
  \item[\textup{(i)}]
  $\kappa_{\{i\}}(G) = \deg_G(i)$.
  \item[\textup{(ii)}]
  For $|S|\ge 2$,
  \[
    \kappa_S(G) = -\,h(G[S]),
  \]
  where $h(G[S])$ is the number of oriented Hamiltonian cycles
  in the induced subgraph $G[S]$, modulo rotation.
  In particular, for $|S|=2$: $\kappa_{\{i,j\}}(G) = -1$ if
  $ij\in E(G)$, and $0$ otherwise, since a Hamiltonian cycle on
  two vertices is a single directed $2$-cycle.
\end{enumerate}
In particular, $\kappa_S(G)\le 0$ for all $|S|\ge 2$.
\end{thm}

\begin{proof}
Since the off-diagonal entries of $L_G$ are $0$ or $-1$, an ordering of $S$ contributes to Theorem~\ref{thm:cycle-expansion} if and only if every consecutive pair is an edge---that is, if and only if it determines an oriented Hamiltonian cycle in $G[S]$.  When it does, the product of $k$ factors of $-1$ gives $(-1)^k$, so $(-1)^{k+1}\cdot(-1)^k = -1$.
\end{proof}


\end{document}
